\clearpage
\section{Discussion}\label{sec:discussion}

\subsection{What the two-land model shows}

The model yields three results relevant to representations of human--environment systems \citep{Schaefer1954,May1973,Meadows1972,Meadows2004,WackernagelRees1996,BranderTaylor1998}. First, conversion can add to total biocapacity even while capital-land area $A_c$ falls; other land flows and quality change determine whether $B$ rises overall. The exact condition for conversion's contribution is in Proposition~\ref{prop:mask}. The technology boost has a fixed upper limit, while debt depends on the size and duration of shortfall (Proposition~\ref{prop:bounded}).

Second, in the reduced fixed-demand comparison ($q=q_{\max}$, $D=0$, no technology boost, and restoration off), the balance relation decreases steadily with $A_c$ until it meets the minimum-area boundary. It does not turn back to form a fold. This finding—and the absence of hysteresis or critical slowing down—applies to that reduced equilibrium relation, not to every trajectory. With one-way conversion and no return flow, $A_c$ remains below its pre-overshoot level after conversion stops, even though it is no longer declining at the balance point (Table~\ref{tab:neg}). Third, recovery requires an explicit restoration flow, and the effect depends on whether its gate opens during a shortfall or a surplus.

SI~\S S7 presents a one-stock model as a separate comparator. Its state $A$ uses comparator-specific units and is not the capital-land area $A_c$. Without debt and with constant parameters, its equilibrium relation is $B(A)=(b+b_G\rho)A-(b_G\rho/A_{\max})A^2$, with $P=B(A)/e$. At the illustrated comparator settings, $b_G\rho<b$, so the relation is monotone; an interior maximum would require different parameter values. The comparator's neutral equilibrium family and selected delay trajectories apply only to that formulation. They do not establish a fold, a universal critical delay, or a collapse mechanism in the two-land model.

These results narrow the conditions under which common early-warning indicators should be expected. Rising autocorrelation and variance are associated with approach to a fold \citep{Scheffer2001}. In this reduced model, drawdown reaches a minimum-area boundary instead, so those indicators need not rise first. Their absence would not show that the system is safe; the relevant signal is the path of $A_c$, which the aggregate alone cannot reveal.

\begin{table}[!htbp]
\centering\footnotesize
\caption{Equilibrium, trajectory, and tested delay-sensitivity properties under the stated assumptions.}
\label{tab:neg}
\begin{tabularx}{\textwidth}{>{\raggedright\arraybackslash}p{0.34\textwidth} l >{\raggedright\arraybackslash}X}
\toprule
Property & Finding & Scope \\
\midrule
Fold on the reduced frozen-demand branch & Absent & $B_{eq}(A_c)$ is strictly monotone; the balance point collides with the typed floor at $\Eceil\approx1.138\,\mathrm{gha}$ (Proposition~\ref{prop:ceiling}) \\
Critical slowing as $E\to\Eceil$ & Absent & The slow eigenvalue remains negative and constant along the reduced branch \\
Hysteresis of the frozen-demand equilibrium relation & Absent & There is at most one feasible balance point; this does not imply history independence of realized states \\
Path dependence of realized $A_c$ & Present, conditional & With $R_{rc}=0$, one-way conversion leaves $A_c$ below its pre-overshoot level after shortfall closure, while $\dot A_c=0$ at equilibrium \\
Discontinuous response to endogenous demand & Not observed & In the reported runs, minimum $A_c$ varies continuously and monotonically with initial population; proximity thresholds depend on tolerance \\
Regeneration-delay sensitivity ($\tau_g$, with $\tau_p=25$\,yr) & Small in tested range & At $A_{c0}=1.0$\,ha, the delayed-quality $10^{-6}$\,ha threshold varies by $0.001$\,cap across $\tau_g=5,20,40$\,yr (SI Table S2); this does not establish invariance outside that range \\
\bottomrule
\end{tabularx}
\end{table}
\clearpage

\subsection{Measurement implications}

Biocapacity in gha is a value-weighted total, not an inventory of its land classes. A $20\%$ rise in the total does not tell us whether all components grew moderately or whether a large gain in one offset losses in others. Proposition~\ref{prop:mask}(iii) gives the formal limit: at fixed per-hectare yields, one total and the conserved-area constraint still leave many feasible combinations of fast, capital, and reserve land.

Related assessment work places $B$ at the linear, perfectly substitutable end of an aggregator family \citep[Lemma~A of][]{Abaee2026b}. An aggregate floor can therefore be met even when a minimum for capital land $A_c$ is not. In that framework, no aggregator with positive substitution elasticity is uniformly safe (Theorem~S2(ii) in \citet{Abaee2026b}). If feasibility must hold throughout a transition, a temporary fall below a land-class minimum remains a violation even if the endpoint later recovers.

For the National Footprint and Biocapacity Accounts, these findings describe account composition, not physical stocks. In Edition 2026, total world biocapacity account value $B$ rises $24.86\%$ from 1961 to 2023, while $B/P$ falls $52.71\%$. The Global Ecological Footprint of Production per person, $E_{\mathrm{prod}}/P$, rises $16.73\%$, and $E_{\mathrm{prod}}/B$ first exceeds one in 1973. Here $B$ and $E_{\mathrm{prod}}$ are annual account amounts in gha, with the annual basis implicit; the per-person quantities are in gha per person, and $E_{\mathrm{prod}}/B$ is dimensionless.

Cropland contributes $+24.06$ percentage points to the total change; Forest biocapacity contributes $-5.95$, grazing land $-0.27$, fishing grounds $-0.42$, and built-up land $+7.44$. The non-cropland account sum is nearly flat because these changes offset one another, including the large rise in built-up land. The same-edition 1961--2022 comparison is reported in SI~\S S5. These are changes in account values, not measurements of forest biomass or condition, soil carbon, wetland extent, fish biomass, or other physical stocks. In particular, the decline in Forest biocapacity does not by itself establish a decline in forest biomass. Built-up land is an NFA accounting class, not ecological capital. Component ratios, contributions, vintage, and endpoint sensitivity characterize account composition but do not establish physical-stock change \citep{Lo2026NFA,Borucke2013,Blomqvist2013,GiampietroSaltelli2014,Galli2016,vdBerghGrazi2015}.

The model result is separate from the NFA account trends. In the frozen-demand baseline (Fig.~\ref{fig:model}a), $A_c$ falls from $1.500$ to $0.807659$\,ha ($46.16\%$), while model biocapacity $B$ rises from $0.594$ to $0.825$\,gha ($39\%$). Although the quality state $q$ rises from $0.9$ to $1.0$, it starts above $q^*=q_{\max}/2$, so $\Ups(q)$ falls from $0.05576$ to $0.05000$\,gha\,ha$^{-1}$ ($10.33\%$). The modeled contribution $\Ups(q)A_c$ falls from $0.08364$ to $0.04038$\,gha ($51.72\%$).

With restoration disabled, one-way conversion leaves $A_c$ below its starting value after the shortfall closes. The area then persists at its lower balance level; it does not continue to decline indefinitely. These are conditional model outputs, not observations or estimates of global physical-stock change.

We did not assemble a common-boundary global physical-stock series for the relevant components, so we report no physical-stock composite or global decline finding. The SI scoreboard is a specification of candidate stock, extent, condition, and annual-flow measures and the requirements for comparing them; it is not an empirical inventory and contains no endpoint values. Until spatial and temporal boundaries, units, uncertainty, overlap, valuation or normalization, and permitted substitution are specified, these heterogeneous measures should remain a nonadditive vector. SEEA Ecosystem Accounting likewise distinguishes ecosystem extent, condition, services, and monetary presentations \citep{UN2021SEEA_EA}.

The external cropland-area and cereal-yield series is an illustrative proxy, not a direct measurement of the model's land classes. Its 1961 starting value is anchored to the NFA cropland share; cereal yield is a sentinel measure, not an all-crops measure, and the area and account boundaries do not match by construction. Here external cropland area $A_f$ is in ha and the index-scaled yield $\hat b_f$ is in gha\,ha$^{-1}$, so $X=\hat b_fA_f$ and the remainder $C=B-X$ are account amounts in gha. That remainder can include boundary and calibration differences, other account components, and conversion-factor changes; it is neither $\Ups(q)A_c$ nor a physical-stock measure. A physical interpretation would require independent, boundary-consistent data on area, condition, yield, and regeneration, plus a validated mapping and valuation. Forest area may inform one part of the effective capital class but cannot represent the whole class or its condition.

\subsection{Policy controls and limits}

The policy analysis compares three kinds of control. A conversion stop triggered at the minimum capital-land area $A_c=\Acmin$ preserves that floor. A ceiling on fast-land area $A_f$ does not guarantee the capital-land minimum unless the stop rule is explicitly linked to it and to the remaining-area limits. A governor that observes only total $B$ cannot guarantee the floor because $B$ does not reveal $A_c$.

The restoration gates also differ: surplus-gated restoration is inactive during shortfall, while a deficit-gated flow acts during shortfall and can reduce provision further. These outcomes connect to the distinction between protective and mobilizing institutional feedback in related work \citep{Abaee2026a}. With a common reserved share $\sigma_f=\sigma_c=\sigma<1$, conversion is inactive while $E\le\sigma B$; with distinct shares, the condition is $E\le\sigma_fY_f+\sigma_cY_c$. This share rule limits the human-available fraction of annual account supply; it is not a spatial map of protected land or an evaluation of a specific conservation proposal. Robustness to disturbances and parameter uncertainty has not been established.

\subsection{Relation to prior work}

The ledger companion develops component-level accounting for material flows, and the assessment companion examines how aggregate measures allow one component to compensate for another \citep{Abaee2026b,Abaee2026c}. This paper adds a delayed land-class model for biocapacity in global hectares. The governance companion studies institutional feedback \citep{Abaee2026a}; the delays here represent ecological regeneration and demographic response, not institutional review.

\subsection{Limitations}

Several limitations govern how to read the model. Ecological-capital land combines diverse assets, so its single quality state, regeneration rule, composition difference, and demand ceiling are effective model quantities—not direct estimates for forests, soils, or wetlands. Most parameters are representative rather than estimated. A cross-section cannot separately identify $b_c$, $b_{G,c}$, and $\rho_c$; recovery records from different systems provide timescale context, not a global calibration of $\rho_c$. They describe regeneration response, not the delay $\tau_g$. The per-person footprint is held constant to isolate the feedbacks among land, biocapacity, and demand; allowing it to change is a further extension.

The composition pattern depends on demand and the mechanisms specified here; under these assumptions, it is bounded. The exact shortfall switch makes stability at $S=0$ one-sided. Reported drawdown thresholds are proximity bands because their values depend on the chosen tolerance and time step. The debt-repayment rate is fixed at $\eta=0.05$\,yr$^{-1}$; the zero-repayment limit, $\eta\to0$, removes the debt equilibrium and is not analyzed. NFA accounts combine area and yield, so they cannot identify those quantities separately. Finally, using current-quality regeneration $G_c(q(t))$ rather than the delayed form $G_c(q(t-\tau_g))$ lengthens the demonstration mask by $1.3$\,yr and shifts the initial-population threshold by at most $0.001$\,cap (SI~\S S3).

\section{Conclusions}\label{sec:conclusions}

Aggregate biocapacity does not reveal how its land classes are changing. In the two-land model, conversion adds to $B$ when fast-land yield exceeds the effective per-hectare value of capital land, even as $A_c$ falls; other terms determine the net change in $B$ (Proposition~\ref{prop:mask}). Under the reduced fixed-demand assumptions, the balance relation has one branch and reaches a minimum-area boundary without a fold, hysteresis, or critical slowing down. The realized path is still history-dependent: with restoration off, $A_c$ remains below its pre-overshoot level after the shortfall closes, but does not continue declining at the balance point. In the representative run, $A_c$ falls $46.16\%$, $\Ups(q)$ falls $10.33\%$, and $\Ups(q)A_c$ falls $51.72\%$ while model $B$ rises $39\%$. These are conditional model results, not estimates of global physical stocks.

In NFA Edition 2026, total world biocapacity account value rises $24.86\%$ from 1961 to 2023, while $B/P$ falls $52.71\%$ and the production footprint per person rises $16.73\%$. Cropland adds $24.06$ percentage points to the account total; Forest biocapacity subtracts $5.95$ points, and built-up land adds $7.44$ as an accounting class, not ecological capital. These account movements are not measurements of biomass, soil condition, wetland extent, fish biomass, or $A_c$. The 1961--2022 sensitivity uses the same account edition and population series. No common-boundary global physical-stock series for the relevant components was assembled, so we report no physical-stock composite or global decline finding. The SI scoreboard is a specification, not an empirical inventory; heterogeneous physical measures remain a nonadditive vector until boundaries, units, uncertainty, overlap, valuation or normalization, and allowed substitution are made explicit. The proxy residual $C=B-\hat b_fA_f$ is neither a physical-capital measure nor the model contribution $\Ups(q)A_c$. A future physical-stock comparison would require independent observations and a validated mapping.

The practical reporting lesson is to publish component accounts alongside the aggregate. The orchard is a metaphor for this point, not evidence of a measured global stock decline. In the analogy, one book records the annual harvest and a second records the trees, soil, and other assets that support later harvests. A real physical-stock claim would require that second book to be filled with independent, comparable observations—not inferred from the account total. A harvest total is not a census of the orchard.

\section*{Acknowledgments}

The author thanks Global Footprint Network for making its data publicly available; Mathis Wackernagel for his thoughtful review; the editors of \emph{Ecological Modelling} for their support; and Arena.ai for access to its online tools.

\section*{Declarations}

\paragraph{Data availability} All empirical inputs are publicly available from the sources cited in the manuscript. Derived analysis data and model code will be deposited in a public GitHub repository and archived in Zenodo; repository and record details will be added after the deposit is complete.

\paragraph{Funding} No funding was received.

\paragraph{Competing interests} The author declares no competing interests.

\paragraph{CRediT authorship contribution statement} \textbf{Amin Abaee:} Conceptualization; Methodology; Software; Formal analysis; Writing -- original draft; Writing -- review \& editing.

\paragraph{Declaration of generative AI in scientific writing} DeepSeek AI and Claude (Anthropic) assisted with computational-reproducibility checks under the author's supervision; the author reviewed and edited the content and takes full responsibility for the publication.
